\documentclass[tikz,border=8pt]{standalone}
\usepackage{tikz}
\usepackage{amsmath}
\usepackage{ctex}
\begin{document}
\begin{tikzpicture}
  % Expectation as weighted average / lever balance point
  % Distribution: p(1)=0.1, p(2)=0.2, p(3)=0.4, p(4)=0.2, p(5)=0.1
  % E[X] = 1*0.1+2*0.2+3*0.4+4*0.2+5*0.1 = 3.0

  \def\Escale{8}  % scale heights

  % Number line (lever)
  \draw[thick] (0.3,0) -- (5.7,0);
  \foreach \x in {1,...,5} {
    \draw[thick] (\x,-0.15) -- (\x,0.15);
    \node[font=\small,below] at (\x,-0.2) {\x};
  }

  % Weights as downward bars + probability labels
  \foreach \x/\p/\h in {
    1/0.1/0.8,
    2/0.2/1.6,
    3/0.4/3.2,
    4/0.2/1.6,
    5/0.1/0.8
  }{
    \fill[blue!60, opacity=0.8] (\x-0.3,-\h) rectangle (\x+0.3,0);
    \draw[blue!80!black] (\x-0.3,-\h) rectangle (\x+0.3,0);
    \node[font=\scriptsize,below] at (\x,-\h-0.12) {$p=\p$};
  }

  % Fulcrum triangle at E[X]=3
  \fill[red!80] (3,-0.15) -- (2.6,-0.7) -- (3.4,-0.7) -- cycle;
  \node[font=\small,red,below] at (3,-0.75) {$E[X]=3$};

  % Balance beam annotation
  \draw[<->,gray,thin] (1,0.4) -- (5,0.4);
  \node[font=\scriptsize,gray,above] at (3,0.4) {加权平均 = 杠杆平衡点};

  % Formula box
  \node[font=\small,fill=yellow!20,draw,rounded corners=2pt,align=center]
    at (3,5.2) {$E[X]=\sum_i x_i\,p(x_i)$\\$=1(0.1)+2(0.2)+3(0.4)+4(0.2)+5(0.1)=3$};

  % Title
  \node[font=\small\bfseries] at (3,6.2) {期望的几何意义：概率加权平衡点};

\end{tikzpicture}
\end{document}
